\subsection{独立表现与新增效用讲述不同故事}
表~\ref{tab:main} 和图~\ref{fig:intervals}a 在相同 spline 探针、对数损失、留出来源和权重下比较检测器独立效用 $U_D$ 与 CDU。MOMENT-FT、MOMENT-ZS 从 $U_D$ 第 4、5 名降至 CDU 第 7、9 名，TranAD 则从第 6 升至第 4。Linear、spline、HGB 下，$U_D$/CDU 排名相关性分别为 0.967、0.667、0.500。基准 VUS-PR 也呈现这种区别：M2N2、TranAD 从 VUS 第 5、6 名变为 CDU 第 2、4 名。

POLY 与两种 MOMENT 的 VUS-PR 接近（0.3832--0.3893），但 POLY 的 spline CDU 比 MOMENT-FT 高 0.00505 比特 [0.00172, 0.00860]，比 MOMENT-ZS 高 0.00593 [0.00241, 0.00961]。三种探针共九项配对比较经最大标准化 Bootstrap 联合校正后，两组区间仍为正。

\subsection{分数可重构性不能决定新增标签效用}
来源留出的 ridge 回归较强地重构 MOMENT-FT 和 MOMENT-ZS 分数（来源宏平均 $R^2=0.799$、0.904；Spearman 为 0.895、0.951），对 M2N2 和 TranAD 则较弱（$R^2=0.054$、$-0.002$）。九模型的重构 $R^2$ 与 spline CDU 的 Spearman 相关性为 $-0.167$（图~\ref{fig:intervals}b）。在这九个检测器中，统计重合与新增标签效用对应较弱：难以重构的分数只有在残余信息改善标签预测时才有价值。

\subsection{条件比较具有可重复性}
Linear、spline 和 HGB 下，SubPCA 与 M2N2 始终排第一、第二，TranAD 与 POLY 组成相同前四名集合（图~\ref{fig:intervals}c）。探针两两排名相关性为 0.717--0.867。Spline 和 HGB 从统计参考中提取的预测信号显著多于 linear logistic：其基底损失分别为 0.2891、0.2869 比特，linear 为 0.3375。线性训练从抽样扩展至全量，基底损失几乎不变（0.337462 对 0.337447），表明差距主要来自探针表达能力，而非训练点数增加。

五个训练抽样种子得到完全相同的 spline 排名。七组参考家族删除的排名相关性为 0.917--1.000；含 7--29 个特征的累积参考保留完整参考的前四名（$\rho=0.933$--1.000）。检测器分数经序列内 z-score 与高斯累积分布函数变换后也保留该领先组（$\rho=0.933$）。最小--最大缩放保留检测器特有的原始分数间距，改变更广泛的排名（$\rho=0.517$），使 POLY 从第 3 降至第 9。这一对照说明采用共同秩接口的理由：主要 CDU 比较衡量异常排序信息，而非检测器特有的分数几何结构。

\subsection{对照验证预期行为}
匹配复制与独立噪声对照均接近零，没有负向对照的区间完全高于零。所有标签相关合成对照明显为正；spline 合成 CDU 为 0.09220 比特 [0.06819, 0.11712]。CDU 区分冗余、无关和互补的标签信号。
